\documentclass[11pt]{article}
\usepackage[margin=25mm]{geometry}
\usepackage{fontspec}
\setmainfont{Latin Modern Roman}
\usepackage{amsmath,amssymb,mathtools}
\usepackage{unicode-math}
\setmathfont{Latin Modern Math}
\usepackage{xurl,hyperref}
\hypersetup{colorlinks=true,urlcolor=blue}
\setcounter{secnumdepth}{0}
\setlength{\parindent}{0pt}
\setlength{\parskip}{5pt}
\setlength{\emergencystretch}{3em}
\newcommand{\sref}[2]{\hyperlink{source:#1:#2}{[S#2]}}
\newcommand{\cataloguescope}[1]{\paragraph{Recorded target.}#1}
\newcommand{\R}{\mathbb R}
\newcommand{\C}{\mathbb C}
\newcommand{\N}{\mathbb N}
\newcommand{\Z}{\mathbb Z}
\newcommand{\Q}{\mathbb Q}
\newcommand{\F}{\mathbb F}
\newcommand{\eps}{\varepsilon}
\begin{document}
% UnsolvedMath ID 61; source catalogue code HIL-006
\setcounter{section}{60}
\section{Hilbert's 6th Problem: Axiomatization of Physics}\label{61}
{\small UnsolvedMath ID 61 · Mathematical Physics}
\subsection{Definitions and mathematical statement}

\paragraph{Shared notation.}$\mathbb N=\{1,2,\ldots\}$; the standard number fields have their usual meanings. All additional conventions are stated locally.


\paragraph{Definitions and notation.}An axiomatization specifies mathematical objects, axioms governing their relations and evolution, and the deductions to be established within that framework. In mechanics, a microscopic model describes individual particles; a continuum model describes fields such as mass density and velocity. A limit from one to the other must specify a scaling, topology of convergence and admissible initial data. Hilbert’s original sixth problem is a programme concerning axioms of physical theories, especially probability and mechanics; it is not a single formal proposition.

\paragraph{Catalogue wording retained for traceability.}
Develop a mathematical framework that axiomatizes physics, particularly mechanics, thermodynamics, and probability theory. Hilbert’s original programme especially calls for an axiomatic treatment of probability, including mean-value methods in mathematical physics, and a rigorous development of limiting processes leading from atomic mechanics to continuum mechanics.
\paragraph{Scope and source qualification.}The catalogue includes thermodynamics as a later framing. Hilbert supplies no fixed final axiom system or universally agreed completion criterion. Existing axiomatizations of particular theories constitute context rather than a settlement of this entire programme.
\subsection{Short English statement}
Give rigorous axioms for physical theories and justify the passage from microscopic particle mechanics to continuum descriptions.
\subsection{Research attempt}
\paragraph{Model, date and investigated scope.}
OpenAI GPT-6 Astra (configured model \texttt{gpt-6-astra}), ultra reasoning; generated 11 October 2026. Proved a free-particle empirical-measure transport limit and exhibited the kinetic moment-closure obstruction. The full physics axiomatization programme is untouched.

\paragraph{Sources and literature check.}
Read the supplied complete definition, catalogue wording, scope qualifications and source list for this canonical ID. The original-source citations are inherited as supplied references; no fresh primary-source or novelty search was performed for this attempt. No current open-status claim is inferred from this calculation.

\paragraph{A precise microscopic-to-continuum benchmark.}Hilbert's programme has no fixed terminal axiom system. To test one component, consider noninteracting particles in \(\mathbb R^d\), with positions and velocities \((x_i,v_i)\) and empirical probability measure \(\mu_N(0)=N^{-1}\sum_i\delta_{(x_i,v_i)}\). Use the distance \(|x-x'|+|v-v'|\) on phase space, and let \(W_1\) be the infimum, over couplings of two probability measures, of the mean distance. Assume finite first moments and \(W_1(\mu_N(0),\mu_0)\to0\). The free flow is \(F_t(x,v)=(x+tv,v)\). For \(0\le t\le T\),
\[d(F_tz,F_tz')\le(1+T)d(z,z'),\qquad
W_1((F_t)_\#\mu_N(0),(F_t)_\#\mu_0)\le(1+T)W_1(\mu_N(0),\mu_0).\]
To prove the second assertion, push any initial coupling forward by \((F_t,F_t)\), integrate the first inequality, then take the infimum. This gives convergence uniformly on every bounded time interval without a stochastic independence assumption.
\paragraph{Continuum equation and a closure failure.}For every compactly supported smooth test function \(\phi\), differentiation under the integral gives \(d\int\phi\,d\mu_t/dt=\int v\cdot\nabla_x\phi\,d\mu_t\), the weak free-transport equation. Yet this kinetic limit is not automatically a closed continuum mechanics system. If a density \(f\) has mass \(\rho=\int f\,dv\) and mean velocity \(u=\rho^{-1}\int vf\,dv\), its momentum equation contains \(\int v\otimes vf\,dv\). At a point, the velocity laws \(\delta_0\) and \((\delta_e+\delta_{-e})/2\), for a unit vector \(e\), have equal mass and mean but different second moments. Therefore mass and momentum alone do not determine the stress.
\paragraph{Gap.}This known free-streaming benchmark establishes one controlled limiting process. It supplies no interaction scaling, collision mechanism, thermodynamic entropy law, or general probability axioms. A closure principle justified from microscopic dynamics is already missing before the full axiomatization programme can be addressed. No universal completion claim is made.

\paragraph{Verification and adversarial audit.}
Checked finite first moments, bounded-time uniformity, test-function differentiation and equal-first-moment counterexample; no interactions or hydrodynamic closure silently assumed. This was an LLM adversarial audit, not independent human verification; no theorem prover was used. Prompts A then B applied to this ID: Checked finite first moments, bounded-time uniformity, test-function differentiation and equal-first-moment counterexample; no interactions or hydrodynamic closure silently assumed. Scope remains limited: Justify an interacting scaling limit and a stress/entropy closure from microscopic hypotheses.

Coordinator checked the pushforward-coupling estimate, finite-first-moment differentiation and unequal second moments with identical mass/mean. This establishes only the free-streaming benchmark. This additional check was model-assisted, not independent human verification.

\paragraph{Tools, process and contributions.}
Python 3 and local shell used to read supplied files and serialize this manuscript; mathematics derived symbolically in the text, with no numerical experiment or formal proof assistant. Contribution prompt A was followed by prompt B. GPT-6 Astra with ultra reasoning drafted the research and performed its stated self-audit; any additional coordinator contribution is identified in the verification and audit. A GPT-6-based Codex coordinator assembled the manuscript and checked repository integration; its exact variant and reasoning configuration were not exposed. No human mathematical contribution, independent human proof check or formal proof certification is supplied. Complete derivations and the consulted source references are included here; private reasoning transcripts are not published.

\paragraph{Final status and first remaining gap.}
\textbf{UNRESOLVED}. Justify an interacting scaling limit and a stress/entropy closure from microscopic hypotheses.

\section{Completion Estimate}
\noindent\textbf{COMPLETION: 0\%}
The calculation is an elementary benchmark or conditional reduction; it establishes no claimed new portion of the complete source target.

\subsection{Sources}
\begin{itemize}
\item[\hypertarget{source:61:1}{S1}] ULAM.AI and the UnsolvedMath Contributors, UnsolvedMath, record 61 (HIL-006). Catalogue curation under CC BY 4.0; definitions and English exposition added locally.
\url{https://huggingface.co/datasets/ulamai/UnsolvedMath}.
\item[\hypertarget{source:61:2}{S2}] David Hilbert, Mathematical Problems (1900), Problem 6, English translation archived by D. E. Joyce.
\url{https://mathcs.clarku.edu/~djoyce/hilbert/problems.html}.
\end{itemize}
\label{end:61}
\end{document}
